\chapter{Environment Setup Reference}
\index{environment setup}\index{AWS CLI}\index{Boto3}\index{cost guardrails}%
This appendix consolidates Chapter~0 into a single reference you can return to before any lab:
account and identity, command-line tooling, SDKs, infrastructure-as-code, and cost guardrails,
followed by a per-chapter table of the tools each week's hands-on project requires. Everything
here assumes the Windows workstation and PowerShell conventions established in Chapter~0.

\begin{notebox}
Use a personal learning account or an organization-approved sandbox for every lab in this book.
Do not run training commands in a production account unless your employer has explicitly approved
the account, Region, budget, and cleanup process.
\end{notebox}

\section{Account and Identity}
\index{root user}\index{IAM Identity Center}\index{multi-factor authentication}%
The account setup sequence is fixed; do it once, in order:

\begin{enumerate}
    \item \textbf{Create the account} with a long-term email address, a non-sensitive account name
    (for example \texttt{data-engineering-sandbox}), a payment card, and a reachable phone number.
    \item \textbf{Enable MFA on the root user} before building anything. Sign out and back in once
    to confirm the second factor is required, then stop using root for daily work.
    \item \textbf{Create a daily administrative identity}, preferring \textbf{IAM Identity Center}
    (workforce sign-in with temporary credentials) over a long-lived IAM admin user. Never create
    access keys for the root user.
    \item \textbf{Set the IAM password policy} (minimum length 14, all character classes) and add
    billing and security contacts.
    \item \textbf{Record the account ID} and store root credentials, MFA recovery notes, and the
    admin profile name in a password manager---never in notes, screenshots, or source files.
\end{enumerate}

\begin{pitfall}
\textbf{Working as root or leaving root without MFA.} An unprotected root sign-in is a high-risk
target: it can change billing, close the account, and lock you out. Root is the owner, not the
operator.
\end{pitfall}

\section{AWS CLI v2 and Profiles}
\index{AWS CLI!profiles}\index{CloudShell}%
Install AWS CLI v2 from the official MSI \cite{awsCliV2UserGuide}, open a fresh PowerShell window,
and verify with \texttt{aws --version}. Then configure access through one of two paths:

\begin{itemize}
    \item \textbf{IAM Identity Center (preferred):} \texttt{aws configure sso} captures the SSO
    start URL and Region; \texttt{aws sso login -{}-profile sandbox} opens the browser flow and
    issues short-lived credentials. No long-lived keys touch the laptop.
    \item \textbf{Access keys (legacy-compatible):} \texttt{aws configure -{}-profile sandbox}
    stores an IAM user's access key under a named profile. Rotate keys periodically and delete
    unused ones.
\end{itemize}

Named profiles keep sandbox and any employer credentials strictly separate; every command in this
book accepts \texttt{-{}-profile <name>}. Set the default Region (\texttt{us-east-1} unless the
chapter says otherwise) in the profile so it travels with the credentials. \textbf{AWS CloudShell}
is the zero-install fallback: a browser shell in the console, pre-authenticated to your console
identity, with the CLI and common SDKs already present.

Verification before any lab:

\begin{lstlisting}[language=bash,caption={Environment smoke test used throughout the book.},label={lst:app-setup-smoke}]
aws sts get-caller-identity --profile sandbox
aws s3 ls --profile sandbox
\end{lstlisting}

The first command proves authentication and shows which account and identity you are acting as;
run it \emph{every} time you switch profiles. The second proves S3 access in the expected Region.

\section{SDKs: Boto3 and Friends}
\index{Boto3}\index{SDK}%
Python scripts in this book use Boto3, the AWS SDK for Python \cite{boto3UserGuide}:

\begin{lstlisting}[language=bash,caption={Installing the Python tooling used across chapters.},label={lst:app-setup-pip}]
pip install boto3 pandas pyarrow redis
\end{lstlisting}

Boto3 reads the same profile chain as the CLI: \texttt{boto3.Session(profile\_name="sandbox")} or
the \texttt{AWS\_PROFILE} environment variable. The recurring pattern is
\texttt{session.client("s3")} for operation-level APIs and
\texttt{session.resource("s3")} where an object-oriented interface reads better. Streaming
chapters add \texttt{redshift-data} (execute SQL without a persistent connection), and the ML
chapters add \texttt{sagemaker} and \texttt{pyspark} where noted.

\section{Infrastructure as Code: CDK and Terraform}
\index{AWS CDK}\index{Terraform}%
Chapters~12 and~24 move from CLI commands to infrastructure as code. Two toolchains appear:

\begin{itemize}
    \item \textbf{AWS CDK} (Chapter~12, orchestration environments): install Node.js, then
    \texttt{npm install -g aws-cdk}, and bootstrap each account/Region pair once with
    \texttt{cdk bootstrap}. CDK synthesizes CloudFormation, so \texttt{cdk synth} is the
    review-before-deploy step \cite{awsCDKDeveloperGuide}.
    \item \textbf{Terraform} (Chapter~24, DataOps): the AWS provider with a remote S3 backend and
    DynamoDB state locking; \texttt{terraform plan} is the review-before-deploy step. CI
    authenticates via OIDC to an IAM role---never stored long-lived keys.
\end{itemize}

Both coexist with the CLI: the CLI is for exploration and smoke tests, IaC is for anything you
intend to keep, repeat, or hand to a teammate.

\section{Cost Guardrails}
\index{AWS Budgets}\index{Cost Explorer}\index{Free Tier}%
Set every guardrail in Chapter~0 before the first lab, and keep them on for the whole program:

\begin{itemize}
    \item \textbf{AWS Budgets:} a monthly cost budget with alerts at 50\%, 80\%, and 100\% of a
    small threshold (for example \$10), plus a forecasted-spend alert.
    \item \textbf{Free Tier usage alerts:} enabled in Billing preferences so approaching an
    allowance emails you before it is exceeded.
    \item \textbf{Cost Explorer:} reviewed weekly, grouped by service, to catch forgotten
    resources---idle NAT gateways, unattached volumes, running clusters, and Serverless workgroups
    left at high base capacity are the recurring offenders.
    \item \textbf{Cleanup habit:} every lab chapter ends with deletion commands; run them. Tag all
    lab resources (\texttt{Environment=Training}, \texttt{Project=<chapter>}) so
    \texttt{aws resourcegroupstaggingapi} can enumerate what a lab created.
\end{itemize}

\begin{pitfall}
\textbf{Assuming Free Tier means no bill.} Free Tier is a discount schedule with limits. Resources
outside the offer, usage above the allowance, idle billable infrastructure, data egress, and
forgotten resources all create charges.
\end{pitfall}

\section{Per-Chapter Tool Requirements}
\index{lab prerequisites}%
The table below distills each chapter's Thursday lab prerequisites. ``CLI'' always means AWS CLI
v2 with a configured profile; ``PS'' means Windows PowerShell. Everything builds on the Chapter~0
sandbox: secured account, named profile, budget alerts.

\begin{center}
\footnotesize
\begin{tabular}{@{}p{1.1cm}p{5.1cm}p{8.6cm}@{}}
\toprule
\textbf{Ch.} & \textbf{Services} & \textbf{Local tooling and prerequisites} \\
\midrule
0 & IAM, STS, S3, Budgets & CLI v2, CloudShell, password manager \\
1 & S3, SQS & CLI; globally unique bucket names; Python + Boto3 for the checker script \\
2 & RDS/Aurora, VPC, Secrets Manager & CLI; small ACU ranges; restrictive security groups \\
3 & DynamoDB & CLI + Boto3; on-demand capacity \\
4 & ElastiCache (Redis), RDS & Python \texttt{redis} client; sandbox VPC with two private subnets; existing RDS/Aurora endpoint \\
5 & Redshift Serverless, S3, IAM & CLI; Python data generator; query editor v2 \\
6 & Redshift Serverless & Chapter~5 workgroup; EXPLAIN for plan inspection \\
7 & Glue crawlers, Athena, Spectrum & Chapter~5 dataset exported as CSV; Athena workgroup with results bucket \\
8 & Lake Formation, Athena & Chapter~7 catalog; two IAM roles (marketing, analytics) \\
9 & EMR (EC2 + Spot fleets), S3 & CLI; PySpark job from Chapter~7 curated zone \\
10 & Glue Studio, Redshift & Chapter~7 catalog tables; Redshift workgroup from Chapter~5 \\
11 & DMS, RDS PostgreSQL, EventBridge & Parameter group with \texttt{rds.logical\_replication=1}; CloudWatch alarms \\
12 & MWAA, Step Functions, SNS & S3 DAG bucket; least-privilege execution role; optional CDK for the environment \\
13 & Kinesis Data Streams & Python + Boto3 producer/consumer \\
14 & Data Firehose, Lambda, Glue table & Chapter~13 stream; Glue table for Parquet conversion \\
15 & Managed Service for Apache Flink & Source and sink Kinesis streams; Studio notebook \\
16 & QuickSight, Redshift, SNS & Chapter~5 warehouse; QuickSight account (free trial tier suffices) \\
17 & Redshift ML, SageMaker, S3 & RedshiftML IAM role; \texttt{ml-artifacts} staging bucket \\
18 & SageMaker Studio, Data Wrangler, Autopilot & Studio domain; governed bucket read through Lake Formation grant \\
19 & SageMaker Pipelines, Feature Store, Registry, Model Monitor & Chapter~18 domain; injected-bad-batch dataset for the gate demo \\
20 & Textract, Comprehend, Lambda, OpenSearch & S3 event notification on \texttt{raw/documents/}; sample scanned PDFs \\
21 & VPC, endpoints, KMS, EC2 + SSM & VPC with no internet gateway; Session Manager for instance access \\
22 & Glue Data Quality, SNS, LF-Tags & Existing Glue job (Chapter~10); deliberately violating batch \\
23 & Macie, EventBridge, Lambda, Object Lock & Object-Locked quarantine bucket; synthetic PII CSV (never real data) \\
24 & Terraform, dbt-redshift, GitHub Actions & GitHub account; S3+DynamoDB backend; OIDC-to-IAM role; Dev/Prod environments \\
\bottomrule
\end{tabular}
\end{center}

\begin{tipbox}
When a lab fails on an unfamiliar error, check the three usual suspects in order: wrong profile
(\texttt{aws sts get-caller-identity}), wrong Region (resources are Regional; the console and CLI
must agree), and a missing Chapter~0 guardrail (budget alerts off, tagging skipped). These three
account for most first-hour failures in every chapter.
\end{tipbox}
